\section{Experimental Design}
\label{sec:experiments}

\begin{intuition}
\textbf{Roadmap: experiments as tests of the law.} The theorem
predicts a regime split: tight clusters admit one-vector summaries,
spread clusters do not. Everything that follows is a test of
that split and its consequences. E1 asks whether the split exists
on a controlled grid. E2/E4/E8 ask whether cluster geometry
organises behaviour across six real corpora and six encoders.
E6 asks whether adaptation between summary families buys anything
beyond picking the right one per cluster. E3 asks how consolidation
trades off against learned quantization across the regime. E5 and
E9 ask whether the answer is stable from $10$K to $1$M vectors and
across query resamplings. E7 asks whether geometry predicts
downstream reader exact-match with a $70$B-parameter Llama reader.
The nine experiments yield $17{,}813$ measured cells ($16{,}000$ of
them from the E1 grid). None of them were designed to crown a
winning method; each was designed to isolate one implication of
the regime split.
\end{intuition}

\subsection{What each experiment is designed to show}
\label{sec:experiments:purpose}

Table~\ref{tab:experiment_matrix} summarises the nine experiments. We
describe them at a high level here; detailed designs appear in the
corresponding results sections.

\begin{table}[h]
\centering
\small
\begin{tabular}{lll}
\toprule
\textbf{Exp.} & \textbf{Question it answers} & \textbf{Cells} \\
\midrule
E1 & Does the tight/spread regime split exist on a controlled grid? & $16{,}000$ \\
E2 & Does cluster geometry organise strategy behaviour across six real corpora? & $503$ \\
E3 & Does geometry, not method family, select between consolidation and quantization? & $338$ \\
E4 & Do full retrieval metrics agree with the identity-coverage story? & $90$ \\
E5 & Is the regime answer stable from $10$K to $1$M vectors? & $23$ \\
E6 & Does adaptation buy anything beyond picking the right summary family per cluster? & $126$ \\
E7 & Does cluster geometry predict downstream reader EM with a $70$B model? & $10$ \\
E8 & Does the geometric prediction reproduce across six sentence encoders? & $198$ \\
E9 & Is the regime answer stable under query resampling? & $525$ \\
\midrule
\textbf{Total} & & $\mathbf{17{,}813}$ \\
\bottomrule
\end{tabular}
\caption{The nine experiments and the part of the law each one tests.
Each row is expanded into its own results section
(\S\ref{sec:e1_validation} through \S\ref{sec:e9}). \textbf{Read
this table as the map of the remaining paper:} every later figure
attacks one of these questions.}
\label{tab:experiment_matrix}
\end{table}

The design has two overlapping guarantees. First, each experiment
answers \emph{one} question about the regime split; the others
either isolate a confounder or extend coverage. Second, each
experiment is independently interpretable: a reader can stop at
the end of any one and take away a complete finding. This matches
the way we ran them in practice: one at a time, with the next
one's design informed by the previous one's result. The order in
which we present them below mirrors the logic rather than the
chronology: E1 establishes the regime split, E2/E4/E8 show that
the split organises real corpora, E6 tests whether adaptation
helps, E3 asks which compression family wins where, and E5/E7/E9
ask whether the answer survives scale, downstream use, and
resampling.

\subsection{The six corpora}
\label{sec:experiments:corpora}

We use six text corpora, chosen to span the $(\deff_{\mathrm{local}},
\dbar)$ plane from the theorem's tight regime to its extreme-coherence
corner.

\begin{itemize}[leftmargin=1.8em, itemsep=2pt]
\item \textbf{drm\_templated} ($n = 4{,}000$, synthetic). Derived from
the Deese--Roediger--McDermott false-memory
paradigm~\cite{roediger1995drm}: $800$ semantically structured facts,
each rendered in five paraphrase templates. Designed as a tight,
low-$\deff$ control: $\deff_{\mathrm{local}} = 2.3$, $\dbar = 0.05$.
This is the \emph{easy} corpus.
\item \textbf{ms\_marco} ($n = 14{,}972$). $1{,}500$ queries from the
MS MARCO passage-ranking benchmark~\cite{nguyen2016msmarco}, each with
roughly ten associated passages. A classic IR benchmark with moderate
per-cluster spread.
\item \textbf{nq\_questions} ($n = 6{,}973$). $500$ groups from
Natural Questions~\cite{kwiatkowski2019nq}, treating each question
group as a cluster. High $\deff_{\mathrm{local}} = 12.6$.
\item \textbf{hotpot\_qa} ($n = 2{,}772$). $1{,}189$ multi-hop question
groups from HotpotQA~\cite{yang2018hotpotqa}. A stress case: tiny
clusters (two or three items), high $\dbar = 0.48$, very small
$\deff_{\mathrm{local}} = 1.5$. The bound is loose here and the
corpus is the hardest in our collection.
\item \textbf{wikipedia\_sections} ($n = 19{,}793$). $379$ sections
from the $20231101$ English Wikipedia dump. High $\dbar$, mid-$\deff$.
\item \textbf{arxiv\_titles} ($n = 11{,}500$). Eleven subject classes
from \texttt{ccdv/arxiv-classification}. Unusual high
$\deff_{\mathrm{local}} = 107.5$: the corpus where the cap-volume
bound goes vacuous first, so we expect learned quantization to
overtake consolidation. It does (\S\ref{sec:e3}).
\end{itemize}

The corpora were selected \emph{before} the experiments ran, with the
criterion that they sample different $(\deff_{\mathrm{local}}, \dbar)$
regimes. They are not tuned for any individual strategy's benefit.

\subsection{The six encoders}
\label{sec:experiments:encoders}

We run the full DRM strategy sweep on six sentence
encoders to separate theorem-predicted universality from encoder-specific
artefacts:
\texttt{BAAI/bge-base-en-v1.5}~\cite{xiao2023bge},
\texttt{BAAI/bge-large-en-v1.5}~\cite{xiao2023bge},
\texttt{sentence-transformers/all-mpnet-base-v2}~\cite{song2020mpnet},
\texttt{sentence-transformers/all-MiniLM-L6-v2}~\cite{wang2020minilm},
\texttt{intfloat/e5-large-v2}~\cite{wang2022e5},
\texttt{nomic-ai/nomic-embed-text-v1.5}~\cite{nussbaum2024nomic}.
Their embedding dimensions range from $384$ (MiniLM) to $1024$ (BGE-large,
E5-large, Nomic), and they differ in contrastive training corpus,
instruction schema, and tokeniser. If the theorem's prediction is a
property of the cluster geometry rather than of any one encoder's
artefacts, the strategy ranking should reproduce across all six
(\S\ref{sec:template_collapse:e8}). It does, with a single exception
(Nomic), which we attribute to the encoder's higher intrinsic $\dbar$
pushing the DRM corpus from tight to spread regime.

\subsection{The five consolidation strategies}
\label{sec:experiments:strategies}

Each experiment evaluates some subset of these five operators.

\begin{itemize}[leftmargin=1.8em, itemsep=2pt]
\item \textbf{Centroid}: $r_k = \mathrm{mean}(\clusterset_k) /
\lVert \mathrm{mean} \rVert$. One representative per cluster;
simplest operator.
\item \textbf{Medoid}: $r_k = \arg\min_{x \in \clusterset_k}
\sum_{y \in \clusterset_k} d(x, y)$. One representative per cluster;
chosen from the cluster itself. Always closer to at least one real
vector than the centroid.
\item \textbf{Importance-weighted}: weighted mean with weights from an
inverse-rank sampler over pairwise similarities. Degenerates to
centroid on isotropic clusters (Corollary~\ref{cor:iw}).
\item \textbf{Selective pruning}: keep the top-$p$\% items by cosine
similarity to the cluster centroid and drop the rest. $p = 50\%$ by
default; we also sweep $p \in \{10\%, 25\%\}$ in E8.
\item \textbf{GAC (Geometry-Aware Consolidation)}: per-cluster router
that picks between centroid, medoid-plus-residual, and pruning
depending on two features of the cluster (\S\ref{sec:gac}). The
production router.
\end{itemize}

Plus a \emph{no-consolidation} reference that keeps all source vectors;
this sets the ceiling for downstream retrieval accuracy.

\subsection{Learned-quantization baselines}
\label{sec:experiments:quantization}

Five non-consolidation families compete against ours on identity
accuracy at matched bytes/vector in E3 (\S\ref{sec:e3}):

\begin{itemize}[leftmargin=1.8em, itemsep=2pt]
\item \textbf{Product Quantization (PQ)}~\cite{jegouPQ2011}: FAISS
implementation, $m \in \{4, 8, 16, 32\}$ chunks, $256$ centroids per
chunk.
\item \textbf{Optimized PQ (OPQ)}~\cite{ge2013optimized}: PQ with a
learned rotation, same chunk schedule.
\item \textbf{LSH}~\cite{charikarLSH2002}: signed random projections to
$\{64, 128, 256, 512\}$ dimensions then reconstructed.
\item \textbf{PCA+int8}: principal components to $\{16, 32, 64, 128\}$
dimensions followed by per-dimension int8 quantisation.
\item \textbf{HNSW-prune}~\cite{malkovHNSW2018}: hnswlib index with
$M = 16$, $ef = 200$, and $\{25\%, 50\%, 75\%\}$ prune-and-rank
retention.
\end{itemize}

All baselines are compared to consolidation on the same $(80\%/20\%)$
train/query split of each corpus.

\subsection{Statistical design: seeds, splits, reporting}
\label{sec:experiments:stats}

Each experiment reports mean performance across three to ten seeds per
cell (exact $n$ per experiment in the respective section). Seeds control
both the random train/query split and any stochastic operator (e.g.\ PQ
initialisation, GAC's random subsample for residual-direction fitting).
Win/tie/loss counts use a tolerance of $10^{-3}$ on mean cap-coverage
error or identity accuracy, chosen to match the one-SD noise floor
across seeds.

All raw results land as line-delimited JSON shards on a persistent Modal
volume, then are reduced to parquet files per experiment. The parquet
files are committed under \texttt{results/} in the repository, along
with the reduction scripts that produced them. Every number cited in
this paper is re-derivable by running \texttt{scripts/make\_figures.py}
and the per-experiment table builders; we verify this end-to-end as part
of Reflection~\ref{sec:limits}.

\subsection{Compute footprint and honesty budget}
\label{sec:experiments:compute}

The full run footprint is $17{,}813$ cells, spread across nine
experiments. Roughly $60\%$ of the wall-clock budget went to E1 (the
$16{,}000$-seed theoretical grid) and $20\%$ to E7 (the
Llama-3.1-70B-Instruct RAG runs on two-H100 nodes with tensor parallelism
2). The remaining $20\%$ covers the six-corpus, six-encoder sweep of
E2--E9. Where we stopped short of the originally planned coverage --- E5
at $1$M points instead of $10$M, E7 at $10$ cells instead of a full
$3 \times 5$ grid, E8 at $198$ cells instead of $1{,}008$ --- we flag the
gap as an honest limitation (\S\ref{sec:limits}), never
as a methodological choice.
